% Generated draft: requires mathematical and source review.
\documentclass[11pt]{article}
\usepackage[margin=25mm]{geometry}
\usepackage{fontspec}
\setmainfont{Latin Modern Roman}
\usepackage{amsmath,amssymb,mathtools,unicode-math}
\setmathfont{Latin Modern Math}
\usepackage{enumitem,needspace,xurl,hyperref}
\setlength{\parindent}{0pt}
\setlength{\parskip}{5pt}
\setlength{\emergencystretch}{3em}
\newcommand{\N}{\mathbb{N}}
\newcommand{\Z}{\mathbb{Z}}
\newcommand{\Q}{\mathbb{Q}}
\newcommand{\R}{\mathbb{R}}
\newcommand{\C}{\mathbb{C}}
\newcommand{\F}{\mathbb{F}}
\newcommand{\A}{\mathbb{A}}
\newcommand{\PP}{\mathbb{P}}
\newcommand{\E}{\mathbb{E}}
\newcommand{\eps}{\varepsilon}
\newcommand{\dd}{\,\mathrm{d}}
\begin{document}
\section*{UnsolvedMath ID: 30002981}
\textbf{Source problem code:} catalogue entry
\subsection{Definitions and mathematical statement}
A nodal curve is a curve whose singularities are ordinary double points. For a linear map \(f\colon U\to W\), its corank is
\[
\operatorname{corank}(f)=\dim\operatorname{coker}(f)=\dim\bigl(W/\operatorname{im}(f)\bigr),
\]
with dimension taken over the coefficient field of the map.

The supplied record denotes a parameter space of nodal curves by \(\mathcal V_{p,n,\delta}\), and associates to a curve \(C\) two maps, the Kemeny--Wahl map \(kw_C\) and the Halic--Wahl map \(hw_C\). It does not supply the construction of this space, the meanings or admissible ranges of \(p,n,\delta\), the ground field, or the domains, codomains, and definitions of these maps. These data remain unresolved; the names alone do not specify a self-contained mathematical problem.

Subject to those missing definitions and hypotheses, the supplied main question is the following: for a component \(Z\) of \(\mathcal V_{p,n,\delta}\) and a general nodal curve \(C\) represented by \(Z\), determine
\[
\operatorname{corank}(kw_C)
\qquad\text{and}\qquad
\operatorname{corank}(hw_C).
\]
The intended parameter ranges and meaning of component must be recovered from the cited source. The word ``general'' is retained from the record; under its usual geometric interpretation, the requested values should hold on a dense Zariski-open subset of the component under consideration.

The original statement additionally asks whether Wahl's conjecture can be extended to nodal curves using these maps. The clean statement omits this clause. Neither the version of Wahl's conjecture nor a proposed nodal analogue is supplied, so this additional question is retained as an unresolved research direction rather than a precisely formulated assertion.
\subsection{Short English statement}
Determine the coranks of the Kemeny--Wahl and Halic--Wahl maps for a general nodal curve in a component of \(\mathcal V_{p,n,\delta}\). The parameter space, parameter ranges, field, and maps require definitions from the source. The original record also asks about extending Wahl's conjecture to nodal curves, without specifying the intended extension.
\subsection{Sources}
\begin{itemize}
\item Local UnsolvedMath registry, supplied record with permanent ID \(30002981\), title ``Coranks of Modified Wahl Maps for Nodal Curves,'' source problem code \texttt{catalogue entry}, in the Oberwolfach Reports Open Problems collection. No individual proposer is identified; the supplied proposed year is 2015.
\item Supplied report citation: \emph{Mini-Workshop: Singular Curves on \(K3\) Surfaces and Hyperkähler Manifolds} (2016), Oberwolfach Reports, pp.\ 2939--2967, DOI 10.4171/owr/2015/51. \url{https://doi.org/10.4171/owr/2015/51}. The citation's year 2016 and the record's proposed year 2015 are retained without reconciliation.
\item Halic, \emph{Modular Properties of Nodal Curves}. \url{https://doi.org/10.1007/s00209-010-0830-2}.
\item Sernesi, \emph{Wahl Map of One-Nodal Curves}. \url{https://arxiv.org/abs/1701.04801}.
\item \emph{The Moduli of Singular Curves on K3 Surfaces}. \url{https://doi.org/10.1016/j.jsc.2015.05.001}.
\item Literature assessment dated \texttt{2026-08-21}: the supplied curation labels the problem ``partially solved'' and reports non-surjectivity of modified Wahl maps in important nodal \(K3\) cases and one-nodal surjectivity in high genus, while stating that coranks for a general component of \(\mathcal V_{p,n,\delta}\) are not determined generally. This is supplied dated curation, not an independently verified assessment. The record itself requests independent verification of time-sensitive status and cited claims.
\item The supplied statement status is \texttt{corrected\_verified}; its accompanying verification note reports comparison with the official Oberwolfach report and preservation of its hypotheses, notation, and mathematical target. This records the supplied verification claim, not independent source consultation for this draft.
\end{itemize}
\end{document}
